\section{Introduction}\label{sec:intro}
Many estimation problems ask for a vector that fits the data well and has
few nonzero entries. A common model gives each continuous variable $x_i$
a binary indicator $z_i$: the variable may be nonzero only if $z_i=1$,
and switching the indicator on costs a penalty $\lambda_i$. With a
quadratic measure of fit, this gives the \emph{indicator quadratic
problem}
\[
 \min\bigl\{x^TQx+c^Tx+\lambda^Tz:\ x_i(1-z_i)=0\ (i\in[n]),\
 z\in\{0,1\}^n\bigr\},
\]
where $Q$ is a symmetric positive-definite matrix. We call a coordinate
with $z_i=1$ \emph{active} and the set of active coordinates the
\emph{support} of a solution; an active coordinate may still be zero.
Once the support is fixed, the best $x$ solves a linear system, so the
difficulty lies in choosing the support. Problems of this form arise in
sparse regression, outlier detection in time series, inference in
Gaussian hidden Markov models, and real-time monitoring of data streams
\cite{BhathenaTrees,BhathenaStructured,GomezBanded}.

Problems of this kind are NP-hard in general (see
\cite{WeiHull,ChoiDiagrams} and the references there). Much work
therefore develops strong convex formulations, often by describing the
closed convex hull of the epigraph of the objective over the feasible
set in an extended space
\cite{AtamturkGomezM,WeiHull,LiuStieltjes,LeeMultiperiod,ChoiDiagrams}.
The same line of work has identified structured cases that can be
solved in polynomial time. These include Stieltjes matrices, whose
off-diagonal entries are nonpositive, with linear coefficients of one
sign \cite{AtamturkGomezM,LiuStieltjes}; matrices whose sparsity pattern
forms a path \cite{LiuPath} or, more generally, a tree
\cite{BhathenaTrees}; matrices with a tridiagonal inverse, which arise
in multi-period models \cite{LeeMultiperiod}; and, through decision
diagrams, quadratics of fixed rank and matrices whose sparsity pattern,
or that of their inverse, is a tree with a fixed number of leaves
\cite{ChoiDiagrams}. Decision diagrams also give an additive
approximation scheme for banded matrices of fixed bandwidth and
condition number \cite{GomezBanded}. For sparse or inverse-sparse
matrices, under conditions on the growth of graph neighborhoods, Choi
et al.\ obtain approximate diagrams whose size is linear in $n$ for each
fixed accuracy \cite{ChoiDiagrams}.

This paper concerns sparse matrices whose \emph{support graph}, with an
edge $ij$ whenever $Q_{ij}\ne0$, has small treewidth. Informally, such a
graph can be assembled from small, overlapping groups of vertices
arranged along a tree, and removing the few vertices shared by two
neighboring groups disconnects the vertices on one side of the tree of
groups from those on the other. Such an arrangement is called a
\emph{tree decomposition}. Trees have treewidth one, and a matrix of
bandwidth $\beta$, with $Q_{ij}=0$ whenever $|i-j|>\beta$, has treewidth
at most $\beta$; Section~\ref{sec:model} gives the precise definitions.
For purely discrete problems on such graphs, the standard exact method
is dynamic programming, which eliminates variables from the leaves of
the tree of groups toward a chosen root. Because the eliminated part
interacts with the rest only through a few shared variables, it suffices
to record its best cost for each combination of their values. With few
shared discrete variables, this table is small.

Indicator quadratics behave differently because the shared variables
include continuous ones. The table then becomes a function: for each
value $t$ of the shared continuous variables (the \emph{boundary
values}), it gives the least cost of the eliminated part with these
variables fixed at $t$. We call such a function a \emph{message}, and an
exact dynamic program needs every message exactly. A message is also
useful in its own right: it gives the least cost and an optimal
completion of the eliminated part for every boundary value at once, so
it can be stored and reused when the rest of the problem changes.
Bhathena et al.\ use this to update solutions online as new observations
arrive \cite{BhathenaTrees}.

Fixing which eliminated coordinates are active leaves a strictly convex
quadratic in the active eliminated variables, and minimizing over them
gives a quadratic function of $t$. A message is therefore the pointwise
minimum of $2^m$ quadratics, one for each support of the $m$ eliminated
variables. Different supports can be optimal at different values of $t$:
even a single eliminated variable with a positive penalty should be
active only for those $t$ that push its best value far enough from zero
to pay for the penalty. A few shared variables thus keep a message
low-dimensional, but they do not limit the number of quadratics needed
to describe it.

On trees, a message depends on one variable, and Bhathena et al.\ showed
that linearly many quadratics suffice \cite{BhathenaTrees}. Beyond
trees, they gave an exact algorithm whose efficiency requires further
conditions on the growth of graph neighborhoods and on how many supports
of each local subproblem are nearly optimal, and they stated that
whether these functions have a compact representation is an open
question \cite[Section~2.2]{BhathenaStructured}. For representations as
minima of quadratics, we show that the answer is negative already at
treewidth two, but that a small random perturbation of the penalties
makes the expected number of quadratics needed polynomial at each fixed
treewidth.

\subsection{Results}
Fix $\sigma>0$ and replace each penalty $\lambda_i$ by
$\lambda_i+\xi_i$, where the \emph{perturbations} $\xi_i$ are
independent uniform draws from a grid of at least $2n$ equally spaced
values in $[-\sigma,\sigma]$. Only the penalties are perturbed. This is
the setting of smoothed analysis, as used for binary optimization
\cite{BeierVocking,RoglinTeng}: an adversary chooses $Q$, $c$, and
$\lambda$, and expectations are taken over the small random
perturbation alone.

Our main result, Theorem~\ref{thm:main}, gives an algorithm that, for
every outcome of the perturbation, solves the perturbed problem exactly
and computes every message of a given tree decomposition exactly on a
box of boundary values with a given radius $R$. Each message is
returned as a \emph{dictionary}: a list of quadratics whose pointwise
minimum equals the message on the box. The dictionary contains the
quadratic of every support that is \emph{optimal somewhere} on the box,
that is, optimal for at least one value $t$ in the box. For each fixed
treewidth, the expected number of bit operations is polynomial in the
input length and in the values, not the bit lengths, of $C$, $H$,
$\mu^{-1}$, $R$, and $\sigma^{-1}$, where $|c_i|\le C$ for all $i$ and
the eigenvalues of $Q$ lie in $[\mu,H]$. The degree of this polynomial
depends on the treewidth.

At a fixed boundary value, the random penalties make near ties
unlikely: the expected number of supports whose cost is within
$\sigma/(2n)$ of the optimum is at most Euler's number $e\approx2.72$.
At every point of a fine finite net of boundary values that covers the
box, the algorithm lists every support within $\sigma/(6n)$ of the
optimum, and only supports within $\sigma/(2n)$, so each list has
expected length at most $e$. For this it needs only an approximate
solver that accepts fixed indicators; for indicator quadratics, a
dynamic program over rounded continuous values provides one. The
quadratics are Lipschitz continuous on the box, so every support that is
optimal somewhere on the box is nearly optimal at a nearby net point and
is listed there; the union of the lists is therefore a dictionary. The
algorithm computes each message directly in this way, from its own
subproblem, instead of passing messages up the tree. The number of net
points and the size of the tables of the rounded dynamic program are
polynomial in $n$, $C$, $H$, $\mu^{-1}$, $R$, and $\sigma^{-1}$, with the
treewidth in the exponent; this is where these values and the treewidth
enter the running time (Section~\ref{sec:main-result}).

Theorem~\ref{thm:dictionary} proves this enumeration step for any family
of supports whose costs are Lipschitz in a parameter and are perturbed
by random penalties. Its ingredients are classical
(Section~\ref{sec:related}); to the best of our knowledge, what is new
is the completeness guarantee for approximate solvers and its use to
cover a whole parameter domain.

The answers are exact for the perturbed instance, not for the original
one. Any exact solution of the perturbed problem gives lower and upper
bounds on the original optimal value that differ by at most $n\sigma$
(Section~\ref{sec:spectral-certificate}). If only an approximate
optimum of the original problem is wanted, the dynamic program over
rounded values gives one directly; the point of Theorem~\ref{thm:main}
is exact messages. A dictionary may contain supports that are never
optimal, and the algorithm does not compute the part of the box on which
each support is optimal.

Two lower bounds explain why the running time depends on numerical
values and why near ties must be controlled, whether by the perturbation
or by conditions such as those of
Bhathena et al.\ \cite{BhathenaStructured}.
Theorem~\ref{thm:hardness} shows that the problem is NP-hard even when
$Q$ has bandwidth two and is arbitrarily close to the identity matrix,
so small treewidth and good conditioning alone do not give a polynomial
exact algorithm unless $\mathrm{P}=\mathrm{NP}$. By scaling the linear
coefficients and penalties of these instances,
Corollary~\ref{cor:magnitude} carries this over to the perturbed problem
with a fixed family of matrices, hence fixed $H$ and $\mu$, and a fixed
$\sigma$. Unless $\mathrm{NP}\subseteq\mathrm{ZPP}$, the polynomial
dependence of Theorem~\ref{thm:main} on $C$ therefore cannot be
replaced by a dependence on the encoding length. Here ZPP is the class
of problems solved by randomized algorithms that never err and run in
expected polynomial time; it appears because the running time of
Theorem~\ref{thm:main} is bounded only in expectation.

Theorem~\ref{thm:fixed-data} shows that messages can be large even when
optimization is easy. Its instances have $2n$ variables, treewidth two,
unit penalties, a matrix close to the identity, and local coefficients
from a fixed finite set. Their global optimum is trivial to find, and
$2^n$ different supports attain it, but each of these supports is the
unique best one for boundary values in a short interval of its own.
Hence one message, a function of a single variable, needs at least
$2^n$ distinct formulas in any exact representation as a minimum of
polynomials or as a piecewise polynomial function. With perturbed
penalties, Theorem~\ref{thm:main} applies to the same instances and
makes the expected dictionary size polynomial. If this message is
needed only within a uniform error $\varepsilon$, a number of its
quadratics bounded independently of $n$ suffices, but this number grows
as a power of $1/\varepsilon$, with matching upper and lower exponents
(Theorem~\ref{thm:pruning}).

Appendix~\ref{sec:representation} concerns representations that depend
only on $Q$ and can therefore be built once and reused for all linear
coefficients and penalties. It studies the inverse-principal polytope,
which Wei et al.\ use to describe the closed convex hull of the
epigraph of $x^TQx$ with indicators \cite{WeiHull}. On a
well-conditioned star, where optimization is easy \cite{BhathenaTrees},
every exact extended formulation (lift) of this polytope by linear,
second-order-cone, or semidefinite constraints with blocks of bounded
order has exponential size (Theorem~\ref{thm:star-face} and
Corollary~\ref{cor:conic}).

The main result is a complexity result, and its polynomial bound has
high degree. We do not propose the algorithm as a practical solver and
report no computational experiments; Section~\ref{sec:discussion}
discusses which parts of the analysis may still be useful in practice.

Section~\ref{sec:model} explains dynamic programming over tree
decompositions for this problem, defines messages and dictionaries,
describes the perturbation, and outlines the proof of the main result.
Section~\ref{sec:related} compares our results with prior work.
Section~\ref{sec:enumeration} proves the general counting and
enumeration results, and Section~\ref{sec:spectral} applies them to
prove Theorem~\ref{thm:main} and to bound the original optimum.
Sections~\ref{sec:hardness} and~\ref{sec:message-limits} prove the two
lower bounds, Section~\ref{sec:discussion} discusses practical relevance
and open questions, and Appendix~\ref{sec:representation} contains the
results on stars.
